% Discussion paragraph: Crossing-right scenario failure mode
% Paste directly into the Discussion section of the paper.

The most pronounced difference between guidance strategies emerged in the \texttt{crossing\_right} scenario, where the no-prediction baseline achieved a perfect $100\%$ success rate while both constant-velocity (CV) and constant-acceleration (CA) prediction methods failed on every episode ($0\%$).

We attribute this to a \emph{forecast-induced virtual-point offset}.  In crossing-right geometries the target approaches from the right with a significant lateral velocity component.  The predictive extrapolators project the target forward by the lookahead horizon ($1\,\text{s}$), shifting the virtual point ahead of the actual target position.  When the projection overestimates the target's turn rate, the pursuer steers too aggressively, overshoots the intercept geometry, and violates the terminal aspect-angle constraint.  When the projection underestimates the turn rate, the virtual point lags behind the actual trajectory, the pursuer turns too late, and the target crosses ahead before intercept can be completed.  The no-prediction method avoids both failure modes by anchoring the virtual point directly on the instantaneous target position, producing a more conservative but robust intercept.

Statistically, the $12.5$ percentage-point overall gap between no-prediction ($75\%$) and prediction-based methods ($62.5\%$) is driven entirely by the ten \texttt{crossing\_right} episodes.  A McNemar-like paired comparison on the four regression scenarios shows that the methods succeed jointly on $30$ episodes, fail jointly on $0$, and diverge on $10$ (binomial $p = 0.0020$).  This confirms that the degradation is \emph{scenario-specific} rather than a general prediction penalty.  Future work will quantify step-level prediction error, virtual-point shift, and mode-switch telemetry to test the overshoot/undershoot hypothesis directly.

\paragraph{CV vs.~CA Equivalence under Constant-Velocity Targets}
Our parametric predictors (CV and CA) produced identical results ($SR = 62.50\%$ each) because all evaluation scenarios used constant-velocity target motion. When target acceleration is zero, the CA model's additional degree of freedom converges to zero, reducing its prediction to the CV form mathematically ($\hat{p}_{CA} = p_0 + v\Delta t + \frac{1}{2}\hat{a}\Delta t^2 = p_0 + v\Delta t = \hat{p}_{CV}$ when $\hat{a} \to 0$). This is not an implementation bug but an expected behavior under the chosen target kinematics. Future work should evaluate on maneuvering targets where CA's acceleration model would provide measurable advantage over CV.
